\section{The X3C checker}
\label{sec:x3c-checker}

\Cref{sec:descent} supplies weight-$h$ binary representatives in
fixed ideal cosets.  For such a representative $\mathbf x$, put
$v_{\mathbf x}:=[\mathbf x]$ and let $\xi_j:=x_{\alpha_j}$.  This section
constructs $U,V\in K$ so that the squared norm of the checker value
$Uv_{\mathbf x}-V$ records the residual
$\mathbf A\boldsymbol\xi-\mathbf1_M$, up to a baseline independent of
$\mathbf x$ and a nonnegative error term.  In
\cref{sec:clean-representatives}, we define clean completions for which this
error term vanishes.  The final corollary records the conclusions used in
\cref{sec:reduction}.

\subsection{Separated offsets}
\label{sec:separated-offsets}

For a selection vector $\boldsymbol\xi\in\{0,1\}^n$ and each $i\in[M]$, define
\[
 \mathcal R_i:=\{j\in[n]:i\in C_j\}
 \qquad\text{and}\qquad
 s_i:=\sum_{j\in\mathcal R_i}\xi_j=(\mathbf A\boldsymbol\xi)_i.
\]
Thus $s_i$ is the number of selected sets containing $i$, and
$\boldsymbol\xi$ is an exact cover exactly when $s_i=1$ for every $i$.  The
quantity encoded by the checker is
\begin{equation}
\label{eq:checker-target-decomposition}
 2\|\mathbf A\boldsymbol\xi-\mathbf1_M\|_2^2
 =
   \underbrace{2\sum_{i=1}^M(s_i^2-s_i)}_{\text{quadratic}}
   -\underbrace{2\sum_{i=1}^M s_i}_{\text{linear}}
   +2M.
\end{equation}

Define the \emph{incidence set}
\[
 \mathcal I:=\{(i,j)\in[M]\times[n]:j\in\mathcal R_i\}.
\]
Every set $C_j$ contains three elements.  We assume here that
$M,n\ge1$ and that every universe element lies in at least one set.  These
conditions are enforced by the preprocessing in
\cref{subsec:prime-and-euclidean}.  Under these assumptions,
\begin{equation}
\label{eq:incidence-count}
 W:=|\mathcal I|=3n,
 \qquad M\le W.
\end{equation}

Let
$0\le\alpha_1<\cdots<\alpha_n=\alpha_\star$ be the Sidon positions from
\cref{lem:sidon}.  The coordinate indexed by $\alpha_j$ will store the
selection bit for $C_j$, so we call $\alpha_j$ a \emph{witness position}.  For
$a,b\in\F_q$, the \emph{offset} from $a$ to $b$ is the cyclic difference
$b-a\in\F_q$.  We separate the offsets associated with different rows of
$\mathbf A$.  Define
\begin{equation}
\label{eq:row-shifts}
 \tau:=2\alpha_\star+1,
 \qquad
 \beta_i:=\tau i\quad(i\in[M]),
 \qquad
 d_{ij}:=\beta_i-\alpha_j\quad((i,j)\in\mathcal I).
\end{equation}
Here $\beta_i$ is the position assigned to row $i$ of $\mathbf A$,
and $d_{ij}$ is the offset from the witness position $\alpha_j$ to
$\beta_i$.  We choose $q$
large enough that
\begin{equation}
\label{eq:row-shift-q-bound}
 q>4\beta_M=4\tau M.
\end{equation}

We use two types of offsets.  The \emph{witness offsets} are the
cyclic differences between distinct witness positions:
\[
\mathcal D_{\mathrm{wit}}
 :=\{\alpha_\ell-\alpha_j:j,\ell\in[n],\ j\ne\ell\}.
\]
By the Sidon property and
$q>4\beta_M>2\alpha_\star$, these differences remain distinct modulo $q$.
The \emph{bad offsets} are the cross-row differences and the sums
of the offsets $d_{ij}$:
\[
 \mathcal D_{\mathrm{bad}}
 :=\{d_{ij}-d_{r\ell}:(i,j),(r,\ell)\in\mathcal I,\ i\ne r\}
 \cup
 \{d_{ij}+d_{r\ell}:(i,j),(r,\ell)\in\mathcal I\}.
\]
In the norm expansion, each witness offset carries a product
$\xi_j\xi_\ell$ from the quadratic term in
\cref{eq:checker-target-decomposition}.  Every correlation at a bad
 offset contributes instead to $\mathcal E(\mathbf x)$.  To preserve the products
at witness offsets while eliminating the correlations at bad offsets, the two
offset sets must be disjoint.  The next lemma proves this separation and bounds the number of
bad offsets.
\begin{lemma}
\label{lem:row-shift-separation}
Under the bound in \cref{eq:row-shift-q-bound}, the $2W$
residues $d_{ij}$ and $-d_{ij}$, for $(i,j)\in\mathcal I$, are nonzero and pairwise
distinct.
Moreover,
\[
 0\notin\mathcal D_{\mathrm{bad}},
 \qquad
 \mathcal D_{\mathrm{bad}}\cap\mathcal D_{\mathrm{wit}}=\varnothing,
 \qquad
 |\mathcal D_{\mathrm{bad}}|=O(W^2)=O(n^2).
\]
\end{lemma}

\begin{proof}
For every $(i,j)\in\mathcal I$,
\[
 \alpha_\star+1=\beta_1-\alpha_\star
 \le d_{ij}\le\beta_M<q/4,
\]
so $d_{ij}$ is nonzero modulo $q$.  Fix distinct
$(i,j),(r,\ell)\in\mathcal I$.

\smallskip
\noindent\emph{Case 1: $i=r$.}
\[
 d_{ij}-d_{i\ell}=\alpha_\ell-\alpha_j\ne0,
 \qquad
 |d_{ij}-d_{i\ell}|\le\alpha_\star.
\]

\smallskip
\noindent\emph{Case 2: $i\ne r$.}
\[
 d_{ij}-d_{r\ell}
 =\beta_i-\beta_r+\alpha_\ell-\alpha_j
 =\tau(i-r)+\alpha_\ell-\alpha_j.
\]
Consequently,
\[
 \alpha_\star+1\le|d_{ij}-d_{r\ell}|
 \le \tau(M-1)+\alpha_\star=\beta_M-\alpha_\star-1.
\]
For all $(i,j),(r,\ell)\in\mathcal I$, one also has
\[
 2\alpha_\star+2=2\beta_1-2\alpha_\star
 \le d_{ij}+d_{r\ell}\le2\beta_M.
\]
Together with $q>4\beta_M$, these bounds place the same-row differences,
cross-row differences, and sums strictly inside $(-q/2,q/2)$, so reduction
modulo $q$ creates no new equalities.  The difference bounds separate residues
of the same sign, and the sum bound separates residues of opposite signs.
Thus the $2W$ residues $d_{ij}$ and $-d_{ij}$ are nonzero and pairwise
distinct.

Every witness offset is represented in
$[-\alpha_\star,\alpha_\star]$.  By contrast, every cross-row difference and
every sum in $\mathcal D_{\mathrm{bad}}$ has a representative in
$(-q/2,q/2)$ with absolute value greater than $\alpha_\star$.  Hence
$0\notin\mathcal D_{\mathrm{bad}}$ and
$\mathcal D_{\mathrm{bad}}\cap\mathcal D_{\mathrm{wit}}=\varnothing$.
The two sets in the definition of $\mathcal D_{\mathrm{bad}}$ contain at most
$O(W^2)=O(n^2)$ offsets.
\end{proof}

\subsection{Checker elements and norm formula}
\label{sec:checker-construction}

Assume throughout this subsection that $q\equiv3\pmod4$.  Fix
$\mathbf x\in\{0,1\}^{\F_q}$ of weight $h$ and put
\[
 v_{\mathbf x}:=[\mathbf x]
 :=\sum_{a\in\F_q}x_a\zeta^a.
\]
Define $\xi_j:=x_{\alpha_j}$ and let
$s_i:=\sum_{j\in\mathcal R_i}\xi_j$ as in \cref{sec:separated-offsets}.  By
\cref{lem:cyclotomic-form} and
\cref{eq:HL-parameters}, $\|v_{\mathbf x}\|^2=H$.

For $d\in\F_q$, define the \emph{cyclic correlation} of $\mathbf x$ at
offset $d$ by
\begin{equation*}
 \mathcal{C}_{\mathbf x}(d):=\sum_{a\in\F_q}x_ax_{a+d}.
\end{equation*}
The indices lie in $\F_q$, so $a+d$ is computed modulo $q$.
Thus $\mathcal C_{\mathbf x}(d)$ counts the ordered pairs
$(a,a+d)\in\F_q^2$ for which $x_a=x_{a+d}=1$.

The preceding decomposition separates the target into quadratic
and linear parts.  We first define $p$, whose shifted monomials encode the
quadratic part through cyclic correlations.  A Gauss sum is then used to
define $U$, so that the quadratic contribution enters
$\|Uv_{\mathbf x}\|^2$ and multiplication by $U$ does not decrease the
canonical norm.  Next, $D$ encodes the linear part, and $V$ combines both
parts in $\|Uv_{\mathbf x}-V\|^2$.
\paragraph{Quadratic term.}
Using the separated offsets $d_{ij}$, define
\begin{equation*}
 p:=\sum_{(i,j)\in\mathcal I}
       \bigl(\zeta^{d_{ij}}+\zeta^{-d_{ij}}\bigr)\in \mathcal{O}_K.
\end{equation*}
By \cref{lem:row-shift-separation}, the displayed $q$-coordinate
expression for $p=\overline p$ has $2W$ terms with pairwise distinct
exponents.  Same-row differences produce
witness offsets, while the other differences and sums produce bad offsets.
The resulting additional correlations are collected in
$\mathcal E(\mathbf x)$ below.

\paragraph{The element $U$.}
Let $\chi$ be the quadratic character of
$\F_q$.  Thus $\chi$ is $1$ on nonzero squares, $-1$ on nonsquares, and $0$ at
zero.  Define the quadratic Gauss sum
\begin{equation*}
 \vartheta:=\sum_{a\in\F_q}\chi(a)\zeta^a\in \mathcal{O}_K.
\end{equation*}

\begin{lemma}
\label{lem:gauss-sum}
The Gauss sum $\vartheta$ satisfies
\[
 \overline\vartheta=-\vartheta,
 \qquad
 \vartheta\overline\vartheta=q,
 \qquad
 \vartheta^2=-q.
\]
\end{lemma}

\begin{proof}
Since $q\equiv3\pmod4$, one has $\chi(-1)=-1$.  Conjugation followed by
$a\mapsto-a$ therefore gives $\overline\vartheta=-\vartheta$.
Since $\chi(0)=0$, writing $a=tb$ in
$\vartheta\overline\vartheta$ gives
\[
 \vartheta\overline\vartheta
 =\sum_{t\in\F_q^\times}\chi(t)
   \sum_{b\in\F_q^\times}\zeta^{(t-1)b}
 =q.
\]
The inner sum is $q-1$ for $t=1$ and $-1$ otherwise.  Combining
$\vartheta\overline\vartheta=q$ with $\overline\vartheta=-\vartheta$ gives
$\vartheta^2=-q$.
\end{proof}

Using \cref{lem:gauss-sum}, define
\begin{equation}
\label{eq:checker-U}
 U:=1+\frac{\vartheta p}{q}\in K.
\end{equation}

\begin{lemma}
\label{lem:checker-U-properties}
The element $U$ has the following properties.
\begin{enumerate}[label=(\roman*)]
 \item $U\overline U=1+p^2/q$.
 \item For every embedding $\sigma:K\hookrightarrow\C$,
 $|\sigma(U)|^2=1+\sigma(p)^2/q\ge1$.  In particular, $U\ne0$.
 \item For every $y\in K$, one has
 $\|Uy\|\ge\|y\|$.
\end{enumerate}
\end{lemma}

\begin{proof}
Since $p=\overline p$, \cref{lem:gauss-sum} gives
$U\overline U=1+p^2/q$, proving~(i).  The value $\sigma(p)$ is real for every
embedding $\sigma$, so applying $\sigma$ to~(i) gives~(ii).  Summing the
lower bound in~(ii) over all embeddings proves~(iii).
\end{proof}

\paragraph{Linear term.}
Define
\begin{equation*}
 D:=-3\sum_{j=1}^n\zeta^{\alpha_j}\in \mathcal{O}_K.
\end{equation*}
Each selected set contributes to three row sums, which motivates
the coefficient $-3$.  Part~(ii) of \cref{lem:checker-U-properties} gives
$U\ne0$, so define
\begin{equation}
\label{eq:checker-V}
 V:=-\frac{D}{q\overline U}
   =-\frac{D}{q-\vartheta p}\in K.
\end{equation}

\paragraph{Norm formula.}
We call $Uv_{\mathbf x}-V\in K$ the \emph{checker value}.  The
correlations in $q^{-1}\|pv_{\mathbf x}\|^2$ that are not determined by the
counts $s_i$ are collected in the following term:
\begin{equation}
\label{eq:checker-additional-correlations}
 \begin{aligned}
  \mathcal E(\mathbf x)
  :={}&
  2\sum_{i=1}^M
   \sum_{\substack{j,\ell\in\mathcal R_i\\j\ne\ell}}
  \bigl(
   \mathcal{C}_{\mathbf x}(\alpha_\ell-\alpha_j)-\xi_j\xi_\ell
  \bigr)
  \\
  &+2\sum_{\substack{i,t\in[M]\\i\ne t}}
   \sum_{\substack{j\in\mathcal R_i\\\ell\in\mathcal R_t}}
   \mathcal{C}_{\mathbf x}
   (\beta_i-\beta_t-\alpha_j+\alpha_\ell)
  \\
  &+2\sum_{i,t\in[M]}
   \sum_{\substack{j\in\mathcal R_i\\\ell\in\mathcal R_t}}
   \mathcal{C}_{\mathbf x}
   (\beta_i+\beta_t-\alpha_j-\alpha_\ell).
 \end{aligned}
\end{equation}
The three lines record, respectively, the additional same-row
correlations, the cross-row difference correlations, and the sum-offset
correlations.

\begin{theorem}
\label{thm:checker-score}
Assume that $q\equiv3\pmod4$.
For every $\mathbf x\in\{0,1\}^{\F_q}$ of weight $h$, define
$\boldsymbol\xi\in\{0,1\}^n$ by $\xi_j:=x_{\alpha_j}$ for $j\in[n]$.  Then
\begin{equation}
\label{eq:checker-score-exact}
 \|Uv_{\mathbf x}-V\|^2
 =B_0+2\|\mathbf A\boldsymbol\xi-\mathbf1_M\|_2^2
   +\mathcal E(\mathbf x),
\end{equation}
where
\begin{equation}
\label{eq:checker-baseline}
 B_0:=H+\|V\|^2+2Wh-2M
       +\frac{2hW-4h^2W^2}{q}
\end{equation}
is a rational quantity determined by the constructed checker and
independent of $\mathbf x$.  Moreover,
$\mathcal E(\mathbf x)\ge0$ for every binary $\mathbf x$.
\end{theorem}

\begin{proof}
\noindent\textbf{Quadratic term.}
For $a\in\F_q$, define
\begin{equation*}
 y_a:=\sum_{(i,j)\in\mathcal I}
      (x_{a-d_{ij}}+x_{a+d_{ij}}).
\end{equation*}
Then $pv_{\mathbf x}=\sum_a y_a\zeta^a$.  Each of the $2W$ shifted copies
of $\mathbf x$ has coefficient sum $h$, so $\sum_a y_a=2Wh$.  For
$r,s\in\F_q$, cyclic reindexing gives
\[
 \sum_{a\in\F_q}x_{a-r}x_{a-s}
 =\mathcal{C}_{\mathbf x}(r-s),
 \qquad
 \mathcal{C}_{\mathbf x}(-r)=\mathcal{C}_{\mathbf x}(r).
\]
For $e=(i,j)$ and $f=(t,\ell)$ in $\mathcal I$, write
$d_e:=d_{ij}$ and $d_f:=d_{t\ell}$.  Expanding $\sum_a y_a^2$ and using
the preceding identities gives
\begin{equation*}
 \sum_{a\in\F_q}y_a^2
 =2Wh+2\sum_{\substack{e,f\in\mathcal I\\e\ne f}}\mathcal{C}_{\mathbf x}(d_e-d_f)
 +2\sum_{e,f\in\mathcal I}\mathcal{C}_{\mathbf x}(d_e+d_f).
\end{equation*}
The diagonal products contribute $2Wh$, and the remaining correlations fall
into three classes.

\smallskip
\noindent\emph{Case 1: $d_e-d_f$ with $e\ne f$ and $i=t$.}
Here $e=(i,j)$ and $f=(i,\ell)$, so $j\ne\ell$ and
$d_{ij}-d_{i\ell}=\alpha_\ell-\alpha_j$.
Since
$\sum_{\substack{j,\ell\in\mathcal R_i\\j\ne\ell}}
\xi_j\xi_\ell=s_i^2-s_i$, this class contributes
$2\sum_i(s_i^2-s_i)$ together with the first line of
\cref{eq:checker-additional-correlations}.

\smallskip
\noindent\emph{Case 2: $d_e-d_f$ with $i\ne t$.}
Here $d_{ij}-d_{t\ell}=\beta_i-\beta_t-\alpha_j+\alpha_\ell$.
These terms form the second line of
\cref{eq:checker-additional-correlations}.

\smallskip
\noindent\emph{Case 3: $d_e+d_f$.}
For all $(i,j),(t,\ell)\in\mathcal I$, one has
$d_{ij}+d_{t\ell}=\beta_i+\beta_t-\alpha_j-\alpha_\ell$.
These terms form the third line of
\cref{eq:checker-additional-correlations}.

Combining the diagonal term and the three cases yields
\begin{equation*}
 \sum_a y_a^2
 =2Wh+2\sum_{i=1}^M(s_i^2-s_i)+\mathcal E(\mathbf x).
\end{equation*}

For $j\ne\ell$, the difference
$\mathcal C_{\mathbf x}(\alpha_\ell-\alpha_j)-\xi_j\xi_\ell$
counts the ordered pairs
$(a,a+\alpha_\ell-\alpha_j)\ne(\alpha_j,\alpha_\ell)$ for which
$x_a=x_{a+\alpha_\ell-\alpha_j}=1$.  The remaining two lines of
\cref{eq:checker-additional-correlations} are sums of correlation
counts.  Thus $\mathcal E(\mathbf x)\ge0$.

Using \cref{lem:cyclotomic-form} and $\sum_a y_a=2Wh$, we obtain
\begin{equation*}
 \|pv_{\mathbf x}\|^2
 =q\left(
 2Wh+2\sum_i(s_i^2-s_i)+\mathcal E(\mathbf x)
 \right)-4W^2h^2.
\end{equation*}

\smallskip
\noindent\textbf{Linear term.}
The coefficient sum of $D$ is $-W$.  Since each selected set contains three
universe elements, $3\sum_j\xi_j=\sum_i s_i$.  Hence
\cref{lem:cyclotomic-form} gives
\begin{equation*}
 \langle v_{\mathbf x},D\rangle
 =-3q\sum_j\xi_j+hW
 =-q\sum_i s_i+hW.
\end{equation*}

\begin{samepage}
\smallskip
\noindent\textbf{Combining the terms.}
Part~(i) of \cref{lem:checker-U-properties}, the definition of
$V$, the equality $\|v_{\mathbf x}\|^2=H$, and the two preceding identities
give
\[
\begin{aligned}
 \|Uv_{\mathbf x}-V\|^2
 &=H+\frac1q\|pv_{\mathbf x}\|^2
   +\frac2q\langle v_{\mathbf x},D\rangle+\|V\|^2\\
 &=H+\|V\|^2+2Wh+\frac{2hW-4h^2W^2}{q}
   +2\sum_i(s_i^2-2s_i)+\mathcal E(\mathbf x)\\
 &=B_0+2\sum_i(s_i-1)^2+\mathcal E(\mathbf x).
\end{aligned}
\]
\end{samepage}
Since $s_i=(\mathbf A\boldsymbol\xi)_i$, the last line is
\cref{eq:checker-score-exact}.  The formula for $B_0$ is independent
of $\mathbf x$, and $B_0\in\Q$ because
$\|V\|^2=\Tr_{K/\Q}(V\overline V)$.
\end{proof}

\subsection{Clean completions}
\label{sec:clean-representatives}

Put $P:=\{\alpha_1,\ldots,\alpha_n\}$.  For each selection vector
$\boldsymbol\xi\in\{0,1\}^n$, define
$\boldsymbol\eta^{\boldsymbol\xi}\in\{0,1\}^P$ by
$\eta^{\boldsymbol\xi}_{\alpha_j}:=\xi_j$.
For $\mathbf u\in\F_q^k$ with $u_0=h$, a \emph{completion} of
$\boldsymbol\xi$ with syndrome $\mathbf u$ is a vector
$\mathbf x\in\mathcal F(\mathbf u;P,\boldsymbol\eta^{\boldsymbol\xi})$.
Thus $\mathbf x$ is binary, has weight $h$ and syndrome $\mathbf u$, and
satisfies $x_{\alpha_j}=\xi_j$ for every $j\in[n]$.

A completion $\mathbf x$ is \emph{clean} if
\begin{equation}
\label{eq:enhanced-clean-correlations}
 \mathcal{C}_{\mathbf x}(\alpha_\ell-\alpha_j)=\xi_j\xi_\ell
 \quad(j,\ell\in[n],\ j\ne\ell),
 \qquad
 \mathcal{C}_{\mathbf x}(d)=0
 \quad(d\in\mathcal D_{\mathrm{bad}}).
\end{equation}
The first equality allows only the ordered pair
$(\alpha_j,\alpha_\ell)$ to contribute at the witness offset
$\alpha_\ell-\alpha_j$.  This ordered pair contributes $\xi_j\xi_\ell$.
The second equality excludes every ordered
pair at a bad offset.  Hence every term in
\cref{eq:checker-additional-correlations} vanishes and
$\mathcal E(\mathbf x)=0$.

The next theorem combines \cref{thm:uniform-completions} with the
fixed-coordinate bound in \cref{cor:fixed-coordinate-ratio} to prove that a
clean completion always exists.
\begin{theorem}
\label{thm:enhanced-clean}
Suppose that $q$ is a sufficiently large prime, $n+2\le T$,
and \mbox{\cref{eq:row-shift-q-bound}} holds.  For every syndrome
$\mathbf u\in\F_q^k$ with $u_0=h$, and every
$\boldsymbol\xi\in\{0,1\}^n$, there is a clean completion
$\mathbf x\in\{0,1\}^{\F_q}$ such that
\[
 \mathbf H_q(k)\mathbf x=\mathbf u,
 \qquad
 \sum_a x_a=h,
 \qquad
 x_{\alpha_j}=\xi_j\quad(j\in[n]).
\]
\end{theorem}

\begin{proof}
Write
$\mathcal F:=\mathcal F(\mathbf u;P,\boldsymbol\eta^{\boldsymbol\xi})$.  By
\cref{thm:uniform-completions}, one has $|\mathcal F|>0$.
For a set $Q\subseteq\F_q\setminus P$ of one or two additional
coordinates, put
\[
 \mathcal F_Q:=\{\mathbf x\in\mathcal F:x_a=1
 \text{ for every }a\in Q\}.
\]
Applying \cref{cor:fixed-coordinate-ratio} gives
\[
 \frac{|\mathcal F_Q|}{|\mathcal F|}
 =O((h/q)^{|Q|}).
\]

\smallskip
\noindent\emph{Case 1: $d=\alpha_\ell-\alpha_j\in\mathcal D_{\mathrm{wit}}$.}
By \cref{lem:sidon,eq:row-shift-q-bound},
$(\alpha_j,\alpha_\ell)$ is the unique ordered pair in $P$ at offset $d$.
This ordered pair contributes $\xi_j\xi_\ell$ to
$\mathcal C_{\mathbf x}(d)$.  For every $a\ne\alpha_j$, the
first equality in \cref{eq:enhanced-clean-correlations} requires
$x_ax_{a+d}=0$, and at least one of $a,a+d$ lies outside $P$.

\smallskip
\noindent\emph{Case 2: $d\in\mathcal D_{\mathrm{bad}}$.}
The second equality in \cref{eq:enhanced-clean-correlations}
requires $x_ax_{a+d}=0$ for every $a\in\F_q$.
By \cref{lem:row-shift-separation}, $d\notin\mathcal D_{\mathrm{wit}}$ and
$d\ne0$.  Thus no ordered pair in $P$ has offset $d$, so at least one of
$a,a+d$ lies outside $P$.

Every required zero product from the two cases has at least one endpoint
outside $P$.  If exactly one endpoint lies outside $P$, no completion
violates the condition when the endpoint in $P$ is fixed to zero.  When
the endpoint in $P$ is fixed to one, the violating completions lie in
$\mathcal F_Q$ with $|Q|=1$ and occupy an
$O(h/q)$ fraction of $\mathcal F$.  If both endpoints lie outside $P$, the
corresponding fraction is $O((h/q)^2)$.

The definition of $\mathcal D_{\mathrm{wit}}$ and
\cref{lem:row-shift-separation} give
$|\mathcal D_{\mathrm{wit}}\cup\mathcal D_{\mathrm{bad}}|=O(n^2)$.  For each
of these offsets, at most $2n$ ordered pairs have exactly one endpoint outside
$P$, and at most $q$ have both endpoints outside $P$.  A union bound shows
that the number of completions violating at least one required zero-product
condition is at most
\[
 O\left(\frac{n^3h}{q}+\frac{n^2h^2}{q}\right)|\mathcal F|.
\]
Since $n\le T=O(q^{1/2000})$ and $h=O(q^{1/1000})$, the relative factor is
$o(1)$.  Therefore some completion in $\mathcal F$ is clean.
\end{proof}

We summarize the checker bounds used in \cref{sec:reduction}.
\begin{corollary}
\label{cor:checker-interface}
For every sufficiently large prime $q\equiv3\pmod4$ satisfying
$n+2\le T$ and the bound in \cref{eq:row-shift-q-bound}, the
following statements hold.
\begin{enumerate}[label=(\roman*)]
 \item Every $\mathbf x\in\{0,1\}^{\F_q}$ of weight $h$, with
 $\xi_j:=x_{\alpha_j}$, satisfies
 \[
  \|Uv_{\mathbf x}-V\|^2
  \ge B_0+2\|\mathbf A\boldsymbol\xi-\mathbf1_M\|_2^2.
 \]
 \item For the center $c$ and syndrome $\mathbf u_c$ defined in
 \cref{sec:descent}, every $\boldsymbol\xi\in\{0,1\}^n$ has a clean completion
 $\mathbf x\in\mathcal F(\mathbf u_c;P,\boldsymbol\eta^{\boldsymbol\xi})$
 such that
 \begin{equation}
 \label{eq:clean-checker-interface}
  v_{\mathbf x}\in c+\mathfrak a_{q,k},
  \qquad \|Uv_{\mathbf x}-V\|^2
  =B_0+2\|\mathbf A\boldsymbol\xi-\mathbf1_M\|_2^2.
 \end{equation}
\end{enumerate}
\end{corollary}

\begin{proof}
Part~(i) follows directly from \cref{thm:checker-score}.
For part~(ii), choose $\mathbf x$ by
\cref{thm:enhanced-clean} with $\mathbf u=\mathbf u_c$.  Cleanliness gives
$\mathcal E(\mathbf x)=0$, so
\cref{eq:checker-score-exact} gives
the norm equality in \cref{eq:clean-checker-interface}.  The vectors
$\mathbf x$ and $\mathbf c$ have syndrome $\mathbf u_c$ and coefficient sum
$h$.  Hence \cref{lem:coset-syndrome} gives
$v_{\mathbf x}=[\mathbf x]\in c+\mathfrak a_{q,k}$.
\end{proof}
